\documentclass[11pt]{article}

\usepackage[a4paper,margin=1in]{geometry}
\usepackage{amsmath}

\title{Research Advancement Plan}
\author{Fred}
\date{}

\begin{document}

\maketitle

\section{Background}
High harmonic generation (HHG) is one of the most important processes
 driving the study of the imteraction of radiation with matter
and will lead to important advances in fundamental science and engineering.
Through HHG, it possible to produce coherent radiations that ouperform the 
external in field in regards to frequency and pulse duration. 

The current theory is based on an approximation that ignores the very fact that 
acceleration of charged particles is in itself, a source of radiation. 
Such radiations as well as produced HHG effect on the whole system is neglected
during the dynamics, among common other simplifications like 
the dipole approximation.
 
\section{Objectives and methods}

Guided by the work of Yangaliev \textit{et al.} [1], 
we want to develop a reliable method of resolution 
of the inhomogeneous Schr\"odinger equation, 
taking into account the radiation emitted by the particle, at various levels.
In other words, we want to build a robust model capable of comparing
results ranging from the classical ansatz 
to the full quantum treatment of the radiation, but also from the 
strong field approximation to the numerical resolution of the 
Time Dependent Schr\"odinger Equation (TDSE), 
in the simplified context of 1D test cases 
before moving on to more realisitic systems.

In order to ensure this transition, we will test our results against, 
some chosen references. Precandidates are suggested down below.
Doing so will allow us to connect our work to the whole litterature, 
still dominated by the work of Lewenstein. 

\section{Timetable}
We shall meet the following deadlines

\subsection{Validation of the main result}
\begin{itemize}
\item \textbf{2026/4/31} \\
\indent Producing the first full quantum compliant results 
		described by Yangaliev \textit{et al} [1], 
        for some emission frequencies only. (not the full spectrum)
         
\item \textbf{2026/5/31} \\
\indent Producing of the full quantum compliant HHG spectra, with 
        converged results in the case of the field strengths 
        $F_0 = 0.03,\mbox{ and }F_0 = 0.1\,\rm a.u.$, for both frequencies 
        $\omega= 0.114\mbox{ and }0.57\,\rm a.u,$
		described by Yangaliev \textit{et al}. [1].

\textbf{Notes}

This will fully validate our model and allow us to reproduce 
the results displayed in the section IV of the reference 
to a satisfactory agreement. Converged results will easily be produced for the 
lesser intensities, but the higher ones will realistically take time 
and will be researched during the subsequent months.


\end{itemize}

\subsection{The adiabatic approximation}
\begin{itemize}
\item \textbf{2026/6/31} \\
\indent In this part, we will study dans comapre our work to 
	Tolstikhin \textit{et al}. [2]. 
	The interest of this reference for us is twofold 
		\begin{itemize}
	\item it will allow us to understand 
		the Adiabatic Approximation (AA) and give us 
		a strong foothold for both the developement of 
		a resolution technic of the TDSE that is not a brute force
		numerical resolution
	\item it will but also will prepare us to our next step
		which is the extension to 
		realistic 3D sytsems.
		\end{itemize}

\textbf{Notes} 

        (N1) The AA will allow us to produce alternative results 
		that may also directly compare to 
		Strong Field Approximation (SFA) 
		results rampant in the litterature for strong external fields.

        (N2) The AA will also allow us to design internal experiments
		for assessing our new fully quantum model in both 
		strong and weak fields.
%        (N2) Validation of our results will be performed by comparing with
%	the work of B. Piraux et al Phys. Rev A 96, 043403 (2017).
%        This choice is motivated by its range, in the tunneling regime
%	when the classical ansatz is expected to be fully valid.
%	Another motivation is its tools, the Sturmian functions 
%	that we would like to revisit. 

\end{itemize}
         
        

\subsection{Extension to realistic systems}

\begin{itemize}
\item \textbf{2026/9/31} \\
    At this stage, we will have a functional fully quantum algorithm of 
		resolution of the TDSE, 
		but also an alternative semi-classical method in the AA. 
		We will aim at extending our work to realistic stystems.
	        Following Tolstikhin \textit{et al}. [3], who already 
		successfully accomplished this by extending [2], to  
		3D systems with finite range potentials, we will have
		comparing points for our results, but most importantly a
		pathway for our approach.


\textbf{Notes} 
    
    (N1) After this point, we will have a two fully functional resolution 
         mechanisms for the TDSE, allowing not only to compare ourselves 
  	 to the general litterature but also to make internal comparison 
	 experiments.
	

\item \textbf{2026/10/31} \\
Last verifications, checks and convergence parameters tuning. If necessary, 
comparison of our results to other works in the litterature 
		for final validation or, following Linblom \textit{et al} [4],
		extension of our work to rotating external fields.

        
\item \textbf{2026/11$\sim$} \\
	Writing a paper and the thesis manuscript. 
\end{itemize}

%\textbf{Tasks:}
%\begin{itemize}
%    \item Analyze the classical/SFA-based Ansatz.
%    \item Compare results with Eq.~(104) when applicable.
%    \item Identify correspondence with the Lewenstein model and SFA.
%\end{itemize}
%
%\textbf{Outcome:}
%Alternative description of HHG with reduced computational cost and conceptual bridge between TDSE-based and semi-classical approaches.
%
%\section*{4. Extension to 3D Systems (by September 30)}
%
%\textbf{Reference:} \\
%K. Ishikawa, ``High-Harmonic Generation,'' in \textit{Advances in Solid State Lasers: Development and Applications} (2010), Sections 2.2 and 2.4.
%
%\textbf{Goal:}
%Extend the current framework to three-dimensional systems including a Coulomb potential.
%
%\textbf{Tasks:}
%\begin{itemize}
%    \item Study standard 3D TDSE formulations (spherical coordinates, partial wave expansion).
%    \item Analyze how the inhomogeneous formulation can be adapted to 3D.
%    \item Evaluate computational requirements and feasibility.
%\end{itemize}
%
%\textbf{Note:}
%Implementation will depend on progress in earlier phases.
%
%\section*{General Remarks}
%
%\begin{itemize}
%    \item Priority is given to obtaining a validated and converged spectrum from Eq.~(103).
%    \item Numerical efficiency will be improved progressively (e.g., reuse of solutions for nearby $\omega$).
%    \item The project aims to connect exact TDSE methods with semi-classical models (SFA, Lewenstein).
%    \item The timeline may be adjusted depending on numerical and theoretical developments.
%\end{itemize}
%
%
\section*{Primary references}
%\begin{itemize}
%%	\item L. Hamonou, T. Morishita,  and O. I. Tolstikhin
%%		Phys. Rev. A 92, 013412 (2012)
%	\item O. I. Tolstikhin, T. Morishita, and S. Watanabe,
%		Phys. Rev. A 81, 033415 (2010)
%	\item O. I. Tolstikhin, and T. Morishita,
%		Phys. Rev. A 86, 043417 (2012)
%	\item M. Ohmi, O. I. Tolstikhin and T. Morishita, 
%		Phys. Rev. A 92, 044302 (2015)
%	\item D. N. Yangaliev, V. P. Krainov and O. I. Tolstikhin, 
%		Phys. Rev A 101, 013410 (2020)
%	\item T. K. Lindblom, O. I. Tolstikhin and T. Morishita, 
%		Phys. Rev. A 104, 023110 (2021)
%\end{itemize}

%	\item L. Hamonou, T. Morishita,  and O. I. Tolstikhin
%		Phys. Rev. A 92, 013412 (2012)
              \noindent	[1] D. N. Yangaliev, V. P. Krainov and O. I. Tolstikhin, 
 		Phys. Rev A 101, 013410 (2020) 

              \noindent [2] O. I. Tolstikhin, T. Morishita, and S. Watanabe, 
	       Phys. Rev. A 81, 033415 (2010) 

	      \noindent [3] O. I. Tolstikhin, and T. Morishita,
               Phys. Rev. A 86, 043417 (2012)

              \noindent	[4] T. K. Lindblom, O. I. Tolstikhin and T. Morishita, 
 		Phys. Rev. A 104, 023110 (2021)

              \noindent [5] M. Ohmi, O. I. Tolstikhin and T. Morishita, 
 		Phys. Rev. A 92, 044302 (2015) \\



\end{document}
